\chapter{The evidence chain behind reliable agents}
\label{closing-evidence-chain}

The most useful result in my retrieval work was a disagreement between two instruments.

Three retrieval measures agreed that the system had become substantially better at placing relevant repository evidence in front of the agent. Across 370 paired tasks, however, the end-to-end reward moved by only +0.0349. Worse, the first interval I reported around that difference was too optimistic: it resampled tasks as though they were independent even though they came from 46 anchor repositories and 20 suites.

The retrieval result was real. I had measured the uncertainty around the downstream result incorrectly.

Understanding that experiment required two parts of this monograph. The stage decomposition from Part IV identified where the improvement occurred. The resampling rules from Part I determined whether the end-to-end difference had been measured adequately at all.

A coding agent's reliability depends on a chain of claims about the system. A strong model, high benchmark score, complete trace, sophisticated retrieval system, or gated deployment cannot establish reliability alone. Each claim gives the next layer something it is allowed to trust.

When an earlier claim is wrong, the failure rarely announces itself downstream. It arrives as a clean score, a confident grader verdict, a plausible trace, a relevant-looking context window, an approved change, or a successfully completed workflow.

The hard problem is deciding which claims about agent behavior deserve to survive contact with the next layer of the system.

\hypertarget{each-layer-determines-what-the-next-may-trust}{%
\section{Each layer determines what the next may trust}\label{each-layer-determines-what-the-next-may-trust}}

Part I sits beneath the rest of this monograph because every later practice is eventually accepted or rejected through a comparison.

One run is one draw from a stochastic system. Before a measured difference can describe the system rather than one favorable draw, the evaluation must characterize local variation, preserve the dependency structure of the sample, and be capable of resolving an effect large enough to matter.

Nineteen practices developed in later parts execute methods introduced in Part I. Removing the explicit evaluation machinery merely leaves those requirements unstated.

Part II turns measurements into operating verdicts.

A model grader, consensus vote, reward model, heuristic score, or proxy metric is an instrument. Each has an operating point, a distribution on which it was calibrated, and errors that can matter asymmetrically. Skalse et al.~(\href{https://arxiv.org/abs/2209.13085}{2022}) showed, in a bounded formal setting over stochastic policies, how restrictive the relationship between a useful reward and an unhacked optimization target can be. The practical lesson is broader than the theorem: a nonconstant proxy should not quietly become final authority merely because it is convenient to optimize.

A grading system therefore produces a verdict under stated conditions. It does not produce truth.

\Needspace{5\baselineskip}

Part III asks whether the record underneath those verdicts is evidence at all.

Four ordinary system failures can corrupt that record in ways no later statistical treatment repairs:

\begin{itemize}
\tightlist
\item
  One identity can reach both primary data and its recovery material.
\item
  A run can die without leaving a durable account of completed work.
\item
  A trace can record tool dispatch without distinguishing whether the external effect committed.
\item
  A quarantine or cleanup policy can remove failed runs before anyone examines them.
\end{itemize}

Even a complete record does not solve attribution automatically. Zhang et al.~(\href{https://arxiv.org/abs/2505.00212}{2025}) gave expert-annotated failure logs from 127 multi-agent systems to the strongest automated attribution methods they evaluated. The best method located the decisive step in only 14.2 percent of cases.

A trace can make an explanation possible, but it cannot manufacture causality.

\Needspace{5\baselineskip}

Part IV determines what evidence reached the model.

A failed implementation should not be assigned to reasoning until the evaluation can answer a more basic sequence of questions:

\begin{itemize}
\tightlist
\item
  Did the required evidence exist in the searchable corpus?
\item
  Was the relevant version present in the index?
\item
  Did retrieval return it?
\item
  Did ranking and context selection preserve it?
\item
  Did it describe the repository state the agent was actually modifying?
\end{itemize}

Stale context is the sharpest example because it can remain perfectly fluent after losing authority. The agent may reason correctly over an interface that no longer exists.

Part V treats human review as a finite system resource rather than an unlimited source of correctness.

Its allocation depends on monitors calibrated in Part II. Its interfaces determine how expensive it is to challenge a claim rather than accept one. Its gates depend on the authority boundaries and evidence preserved in Part III.

A gate that cannot change the execution path records assent. It does not provide control.

Part VI depends on artifacts produced by every earlier part. Routing and scheduling policies should be judged by replaying recorded arrivals, holding demand fixed, and forcing sophisticated policies to compete against cheap alternatives.

That requirement appears repeatedly in the literature because complexity has a habit of beating an absent baseline. Kapoor et al.~(\href{https://arxiv.org/abs/2407.01502}{2024}) found that simple model retries could match more elaborate agent architectures on a function-level coding benchmark at a fraction of their inference cost. Chen et al.~(\href{https://arxiv.org/abs/1608.07617}{2016}) made the same methodological point in search-based software engineering: an optimization method should first beat inexpensive random sampling.

The exact baseline changes, but the obligation does not.

The dependency chain also explains why some repairs are so expensive. Later machinery cannot reliably compensate for an earlier defect:

\begin{itemize}
\tightlist
\item
  More samples do not repair a task distribution that excludes production work.
\item
  More judges do not repair a rubric experts interpret inconsistently.
\item
  More agents do not repair a retrieval boundary that treats an empty result as authoritative.
\item
  More context does not repair a run that has lost its governing specification.
\item
  Better scheduling does not repair a work ledger that cannot identify which attempt owns the right to publish.
\item
  More observability does not repair an external effect whose identity was never recorded.
\end{itemize}

In each case, the added machinery is evaluated through the instrument that was already defective.

This repair asymmetry lets problems propagate forward cheaply while making downstream correction expensive.

\hypertarget{engineering-controls-can-precede-prevalence-evidence}{%
\section{Engineering controls can precede prevalence evidence}\label{engineering-controls-can-precede-prevalence-evidence}}

The evidence in this monograph is uneven, and that unevenness does not track operational consequence.

Strong-graded items make up 56 percent of the evidence in Part I but only 10 percent in Part III. Chapter 8 contains no strong-graded or direct scholarly evidence item even though its containment practices may be among the first controls needed by a system with production write access.

Within the 52-practice consequence-ranking subset, Spearman's rho between urgency rank and the presence of at least one strong evidence item is -0.004. In that selection, evidence strength and operational urgency are effectively unrelated.

Engineering sometimes has to move before prevalence has been estimated.

\texttt{Contain agent blast radius} ranked first by consequence in the catalog, yet its incident support consists of two corroborating practitioner reports: one involving a production workflow and another involving database deletion. Those reports establish possibility. They do not establish frequency.

I still recommend the control because its mechanism can be inspected directly. If one credential can alter both production state and the material required to recover that state, those resources share a failure domain. The authority boundary can be changed. A prohibited write can be exercised against a test resource. Ordinary and recovery authority can be separated. The control is reversible and its local effect is observable.

This is a mechanism-based justification, not a claim that the literature has proved a universal effect size.

Some practices in this monograph are supported by measured effects. Others are supported by a structural failure argument:

\begin{enumerate}
\def\labelenumi{\arabic{enumi}.}
\tightlist
\item
  The failure is possible under the current authority or state model.
\item
  Its consequence is material.
\item
  The proposed control changes the relevant mechanism rather than a proxy for it.
\item
  The change is reversible or bounded.
\item
  The resulting boundary can be exercised directly.
\end{enumerate}

The following are observations of local protections, not estimates of production frequency:

\begin{itemize}
\tightlist
\item
  A denied write.
\item
  A durable completion record that survives a worker crash.
\item
  A stale completion rejected by a newer ownership generation.
\item
  A replay in which the same recorded arrivals face two schedulers.
\end{itemize}

Each observation shows whether a claimed local protection exists.

The imperative headings in thin-evidence chapters should be read in this sense. They specify an engineering action and the observation that can falsify its local justification. They do not imply a universal prevalence estimate or effect size.

A mechanism that cannot be reduced to an executable observation has not yet been specified well enough to rely on.

Important gaps remain. Two studies would materially strengthen the record developed here:

\begin{itemize}
\tightlist
\item
  A controlled comparison of isolation and authority designs scored against observed incidents.
\item
  A recovery benchmark with a published fault menu, named kill placements, recurring failures, lost acknowledgements, and upgrades across runtime versions.
\end{itemize}

Part VI remains more transfer-heavy still. Its topology, admission-control, partitioning, and scheduling proposals are executable research questions for coding-agent fleets. They should not be mistaken for established production effects.

Useful engineering hypotheses need not masquerade as settled science. They need claims that state what kind of evidence supports them.

\hypertarget{start-smaller-than-this-book}{%
\section{Start smaller than this book}\label{start-smaller-than-this-book}}

A production system does not need every mechanism described in these chapters before it can improve.

It does need enough instrumentation to tell whether an improvement occurred and enough durable state to show what happened when something failed.

The literature does not provide universal values for repeat count, inter-rater agreement, context limits, refresh cadence, routing cost, retry count, or re-solve frequency. The values below are starting points for a new system, not transferred effect estimates.

\begin{longtable}[]{@{}
  >{\raggedright\arraybackslash}p{(\columnwidth - 4\tabcolsep) * \real{0.22}}
  >{\raggedright\arraybackslash}p{(\columnwidth - 4\tabcolsep) * \real{0.43}}
  >{\raggedright\arraybackslash}p{(\columnwidth - 4\tabcolsep) * \real{0.35}}@{}}
\caption{Personal starting points for a new system. These values are author defaults, not transferred effect estimates.}\tabularnewline
\toprule\noalign{}
\begin{minipage}[b]{\linewidth}\raggedright
Decision
\end{minipage} & \begin{minipage}[b]{\linewidth}\raggedright
Personal starting point
\end{minipage} & \begin{minipage}[b]{\linewidth}\raggedright
What changes it
\end{minipage} \\
\midrule\noalign{}
\endfirsthead
\endhead
\bottomrule\noalign{}
\endlastfoot
\textbf{Repeated evaluation} & Use \(k=3\) paired runs per item for inexpensive screening, but do not promote from that screen alone. Size release comparisons from a pilot and the smallest decision-relevant effect. & Higher variance, clustered tasks, or smaller meaningful effects require more runs. \\
\textbf{Reliability reporting} & Report both \(\mathrm{pass}@k\), where any attempt succeeds, and \(\mathrm{pass}^{k}\), where every attempt succeeds, with \(k\) printed beside the result. Begin with \(k=3\). & The deployed retry policy and the cost of intermittent failure determine the useful \(k\). \\
\textbf{Human-label agreement} & Treat Cohen's or Fleiss's kappa below 0.60 as a rubric-debugging trigger. A value above 0.60 is not a correctness gate. Promotion still depends on class-specific errors against expert adjudication. & Class imbalance, ambiguous labels, and high-consequence mistakes make the confusion matrix more informative than kappa alone. \\
\textbf{Promotion threshold} & Before execution, require either a three-percentage-point absolute success gain at no material cost increase or a ten-percent cost reduction at no success loss. Treat smaller effects as uncredited until a decision justifies them. & Task value, baseline success rate, and operating cost should replace these round numbers. \\
\textbf{Failure review} & Read twenty stratified failures after a material model, evaluation apparatus, or policy change. Preserve an abstain category when the trace cannot support attribution. & Rare high-consequence failure classes should be oversampled regardless of frequency. \\
\textbf{Context restart} & Restart or consolidate at 60 percent of the shortest context length at which the weakest required task shows material degradation in a local sweep. & New models, tools, instruction files, or task mixtures require another sweep. \\
\textbf{Freshness} & Require exact repository-revision identity for code-bearing retrieval and reject an index whose source revision is unknown. & A coarser content-addressed equivalence is acceptable only when the indexer can prove it. \\
\end{longtable}

The apparatus itself belongs in the cost model.

Each serious comparison should record inference and tool cost, elapsed time, labeling and adjudication hours, reviewer queue time, storage, and maintenance work. The denominator should be accepted work rather than model calls.

A control whose measurement cost remains invisible will eventually be bypassed, removed, or defended with stale evidence.

When capacity is constrained, remove complexity from the outside inward. Cut dynamic scheduling and learned routing before cutting the records needed to evaluate them. Cut multi-agent debate, specialist roles, and elaborate aggregation before cutting the acceptance check. Cut redundant model judges and retrieval lanes before cutting the ability to reconstruct which artifact was accepted. Reduce benchmark breadth before abandoning a small stratified set of consequential tasks.

Rich interfaces and secondary taxonomies can also go if their questions remain answerable from the durable record.

Four things should survive almost to the floor:

\begin{enumerate}
\def\labelenumi{\arabic{enumi}.}
\item
  \textbf{An executable acceptance check.}
\item
  \textbf{Separation between ordinary authority and recovery authority.}
\item
  \textbf{Versioned identity for the task, input state, configuration, attempt, and accepted artifact.}
\item
  \textbf{A paired comparison against the cheapest credible baseline.}
\end{enumerate}

A small team can operate those controls without reproducing the full apparatus described in this monograph.

This minimum viable reliability system can show whether the work passed and which state and configuration produced it. It can keep a failed worker from silently retaining authority and show whether a more complicated system beat the simpler one.

Everything added after that should have to earn its place.

\hypertarget{some-uncertainty-should-remain-visible}{%
\section{Some uncertainty should remain visible}\label{some-uncertainty-should-remain-visible}}

Several questions in this book remain research problems rather than missing configuration defaults.

No controlled result in the reviewed set establishes that a dynamic dependency graph outperforms a well-designed fixed schedule for repository work. Propagation of deletion through summaries, embeddings, caches, and graph edges remains poorly measured. The optimal recovery architecture for long-running coding agents is not settled. Neither are universal grader thresholds, context policies, retry counts, or fleet scheduling rules.

Many empirical conclusions here are also dated snapshots. Models, context windows, retrieval systems, agent runtimes, and costs will change. Benchmarks will become contaminated or obsolete.

The structural claims are harder to age out:

\begin{itemize}
\tightlist
\item
  One stochastic run is one draw.
\item
  A component that did not execute did not cause the observed result.
\item
  An instrument has to be calibrated for the decision it is being asked to make.
\item
  A trace can preserve evidence without proving causality.
\item
  An empty retrieval result is not proof that relevant evidence does not exist.
\item
  A stale artifact can remain fluent after losing authority.
\item
  An acknowledged workflow transition does not prove an external effect occurred exactly once.
\item
  A worker that has lost ownership must also lose the ability to make consequential mutations.
\item
  One identity that reaches production data and recovery material defines one failure domain.
\item
  Additional machinery should beat the cheapest credible alternative under the same workload before it earns adoption.
\end{itemize}

Those claims describe constraints on reasoning about systems, not preferences for today's agent stack. They are the part of this book I expect to last longest.

\hypertarget{the-system-around-the-model}{%
\section{The system around the model}\label{the-system-around-the-model}}

The easiest way to misunderstand a coding agent is to look only at the model.

A production result is jointly determined by the task presented to the system, the version of the world the system can see, the evidence retrieval makes available, the state retained across steps, the tools and permissions through which actions occur, the workflow that decides what may happen next, the acceptance path that decides what counts, and the evaluation that tells us whether any of those choices helped.

The model matters enormously, but it is only one component in the causal path.

A better model can fail inside a weak execution system; better retrieval can disappear inside a noisy end-to-end score; a complete trace can support the wrong attribution; a successful retry can duplicate an external effect; and an elaborate fleet scheduler can optimize a workload whose ownership semantics were never correct.

The practical consequence is a debugging order:

\begin{enumerate}
\def\labelenumi{\arabic{enumi}.}
\item
  Begin with the comparison you currently trust most.
\item
  Preserve the per-item result.
\item
  Ask whether repeated runs support the difference.
\item
  Check whether the verdict instrument agrees with expert adjudication on the errors that matter.
\item
  Confirm that the durable record can reconstruct the accepted attempt and artifact.
\item
  Verify that the evidence required by the task existed, was current, and reached the model.
\item
  Exercise the consequential boundary under failure.
\item
  Add machinery only after those checks pass.
\end{enumerate}

The six-step minimum pass is the shortest practical route through this dependency chain. It exposes the claims a system already depends on without prescribing a single agent architecture.

A reliable agent system preserves the evidence needed to understand a failure, prevents stale authority from surviving recovery, keeps uncertain measurements uncertain, and requires added complexity to improve accepted work.

The model can propose the work. The surrounding system determines what evidence the proposal used, what it could change, whether the result was accepted, whether the effect occurred, and what evidence survived.

\section*{Sources and evidence}

No evidence is introduced here. Each identifier below is carried by the chapter named, and the evidence grouping comes from that chapter's record.

\begin{itemize}
\tightlist
\item
  Strong evidence: Skalse, Howe, Krasheninnikov \& Krueger (2022). \emph{Defining and Characterizing Reward Hacking}. NeurIPS 2022. arXiv:2209.13085. Chapter 6, \texttt{layer-signals-beyond-single-proxy}.
\item
  Strong evidence: Zhang, S., et al.~(2025). \emph{Which Agent Causes Task Failures and When? On Automated Failure Attribution of LLM Multi-Agent Systems}. ICML 2025. arXiv:2505.00212. Chapter 11, \texttt{keep-humans-in-failure-attribution}.
\item
  Directional evidence: Kapoor, Stroebl, Siegel, Nadgir \& Narayanan (2024). \emph{AI Agents That Matter}. arXiv:2407.01502. Chapter 2, \texttt{report-cost-accuracy-pareto}.
\item
  Strong evidence: Chen, J., et al.~(2016). \emph{Sampling as a Baseline Optimizer for Search-Based Software Engineering}. arXiv:1608.07617. Chapter 19, \texttt{replay-traces-before-policy-changes}.
\end{itemize}

The part-level and chapter-level evidence shares restated above were recomputed from the companion catalog: Part I has 18 of 32 items grouped as strong, Part III has 3 of 30, Part IV has 20 of 42, and Part VI has 8 of 29. Chapter 8 has no strong or direct scholarly evidence item.
